This manuscript studies the Riemann Hypothesis through a proposed balance-functional perspective built on the reflection symmetry of the completed zeta function $\xi(s)=\xi(1-s)$ and the direct object $\Phi(s):=\Gamma(s)(1-s)\zeta(s)$. The framework organizes a critical-line strategy around a symmetry-and-balance principle: the nontrivial zero set is modeled as constrained by a balance condition coupled to completed-function symmetries, with the intended conclusion that sustained balance occurs only at $\RePart(s)=\tfrac12$. Earlier drafts used spectral-coherence language for related intuition; the current manuscript centers the balance-functional formulation. The manuscript is presented as a proposed proof framework under review and as material for mathematical audit, with formal targets, assumptions, and dependencies stated explicitly. Its central unresolved burden is the off-line exclusion bridge: proving, under explicit hypotheses, that any zero with $\RePart(s)\neq\tfrac12$ forces a failure of the proposed balance condition.
